% Propietario conservado: /Users/ruben/Documents/New project/output/REVISION_ARGUMENTAL_APERTURAS_CIERRES_20260915/VII/delivery/MANUSCRIPT_VII_ES_EN_COMPLETE/source_es/manuscrito/33_corriente_espinorial_y_densidad.tex
% Artículo VII. Incorporación aditiva autorizada, 10-09-2026.
% RESULTADO_RECUPERADO: representación real de Clifford y CAR:
% TEORIA_HOLOGRAFICA_INTEGRAL_MASA_HMT_MD_REV2_2026-08-20/fuente/modulos/
% 18_cuantica_estadistica_y_termodinamica.md, líneas 131--181 y 218--255;
% acción de primer orden: 14_topologia_exoticos_y_gravedad.md, 459--517;
% distinción corriente media/segundo momento y dilución:
% HMT_OBRAS/OBRA_II/chapters/11_cartan_bianchi_cosmologia.tex, 201--230.
% PDF REV2: b56bb38d10b9a4f4d0de47b7da8010fb31a6e8e8684a3bfd4237e37d7a3dddfa.
% FORMALIZACION_NUEVA: acción espinorial afín explícita, contracción
% normalizada mediante 31, observable cuártico CAR y funcional de amplitud.
% CERTIFICADO_NUEVO: cálculo racional de la forma cuadrática en las cuatro
% componentes de una trivector y evaluación de la realización CAR finita.
% Convención explicitada: las matrices g tienen firma (-+++); A_D=-i g^0
% y la acción cinética simétrica lleva hbar/2. La combinación equivalente
% i hbar Gamma^I usa Gamma^I=-i g^I y firma opuesta. No se mezclan ambas.
% Dependencias: cotetrada/conexión de 30d; constantes internas heredadas;
% corriente e inversas de 31. Estado material, celda y normal orden son
% datos de la realización declarada, no constantes universales extraídas.
% La acción afín de este fragmento y la acción de extensión mínima de 30e
% conservan sus problemas variacionales propios. No se identifican sus
% corrientes. No se cambia 50: su dominio de amplitud positiva se concreta.

\section{Corriente espinorial y amplitud torsional del estado material}
\label{vii:vii:sec:corriente-espinorial-densidad}

El desarrollo precedente y la realización espinorial que sigue articulan
dos problemas variacionales con antecedentes geométricos comunes. La
acción de extensión mínima \(S_{\min}=\mathfrak a_m E\), construida en
la sección~\ref{vii:vii:sec:corriente-extension-minima}, incorpora la
dependencia de la incidencia transportada y de su mínimo respecto de la
conexión; su eliminación torsional conserva esa dependencia completa en
la sección~\ref{vii:vii:sec:eliminacion-torsional-no-lineal}. La acción
espinorial \(S_D\) de \eqref{vii:vii:eq:accion-dirac-material} acopla la
misma conexión a los campos de su representación y es afín en la
contorsión al fijar la cotetrada y los campos. Cada corriente se obtiene
variando su propia acción. Una preparación minimizante puede determinar
el estado material inicial de la realización espinorial; la sustitución
de campos dependientes de la conexión conserva, además, los términos
del diferencial compuesto. Esta distinción sitúa el cálculo de amplitud
siguiente en su acción y en su dominio variacional específicos.

La representación espinorial del transporte y la respuesta algebraica
de Cartan--Holst permiten calcular la intensidad torsional como un
funcional del estado material. Su evaluación conserva tres datos:
la representación de los campos, el segundo momento de su corriente y
el volumen respecto del cual se define la densidad. El coeficiente de
acoplamiento resulta de la acción y de su inversión; el estado inicial
determina la amplitud material.

\subsection{Representación espinorial y acción material afín}
\label{vii:vii:subsec:accion-espinorial-afin}

Conservamos la cotetrada, la orientación y los acoplamientos de la
sección~\ref{vii:vii:sec:corriente-respuesta-cartan}. La realización
espinorial utiliza el módulo complejo de dimensión cuatro obtenido al
complejificar las matrices reales
\begin{equation}
 \begin{gathered}
 J=\begin{pmatrix}0&-1\\1&0\end{pmatrix},\qquad
 R=\begin{pmatrix}0&1\\1&0\end{pmatrix},\qquad
 Z=\begin{pmatrix}1&0\\0&-1\end{pmatrix},\\
 g^0=J\otimes I_2,\qquad g^1=R\otimes I_2,\qquad
 g^2=Z\otimes R,\qquad g^3=Z\otimes Z.
 \end{gathered}
 \label{vii:vii:eq:clifford-material-explicito}
\end{equation}
La multiplicación da
\(\{g^I,g^J\}=2\eta^{IJ}I_4\), con
\(\eta=\operatorname{diag}(-1,1,1,1)\). Las matrices se distinguen
del parámetro escalar \(\gamma\) del operador de Holst. Definimos
\begin{equation}
 g_I=\eta_{IJ}g^J,\quad
 \mathbb B_{IJ}=\tfrac14[g_I,g_J],\quad
 A_D=-i g^0,\quad \bar\psi=\psi^\dagger A_D,\quad
 \Omega=\tfrac12\omega^{IJ}\mathbb B_{IJ}.
 \label{vii:vii:eq:adjunto-conexion-espinorial}
\end{equation}
El producto de Clifford verifica
\[
 [\mathbb B_{IJ},g^K]=\delta_J^K g_I-\delta_I^K g_J.
\]
Así, la misma conexión métrica actúa en los campos mediante
\(D\psi=\mathrm d\psi+\Omega\psi\) y
\(D\bar\psi=\mathrm d\bar\psi-\bar\psi\Omega\).
El transporte sobre una ruta se obtiene integrando esta conexión en
su representación espinorial; conserva la conexión y la ruta de la
realización geométrica precedente.

Sea \(\eta_I=\iota_{E_I}\operatorname{vol}_e\), de modo que
\(e^J\wedge\eta_I=\delta_I^J\operatorname{vol}_e\). Para un
campo de tipo dimensional \([\psi]=L^{-3/2}\), tomamos la acción
material mínima
\begin{equation}
 S_D=\int_M\left\{
 \frac{\hbar}{2}
 \bigl[\bar\psi g^I D\psi-(D\bar\psi)g^I\psi\bigr]\wedge\eta_I
 -mc\,\bar\psi\psi\operatorname{vol}_e\right\}.
 \label{vii:vii:eq:accion-dirac-material}
\end{equation}
Aquí \(\hbar=\hbar_{\rm int}\), \(c=c_{\rm int}\) y la masa
\(m\) es la lectura de la especie material en la carta elegida.
Estos valores son salidas anteriores a esta realización. La matriz
\(A_D\) es hermítica y \(A_Dg^I\) es antihermítica. Por ello el
término cinético simétrico tiene la convención real correspondiente a
\eqref{vii:vii:eq:clifford-material-explicito}. En esta firma, la ecuación
sin torsión tiene término principal \(\hbar g^I D_I\); la notación
equivalente \(i\hbar\Gamma^I D_I\), con \(\Gamma^I=-i g^I\),
emplea la relación de Clifford de firma opuesta. Esta precisión fija
conjuntamente adjunto, firma y factor imaginario de la acción.

Los campos materiales y la cotetrada permanecen fijos en la variación
de la conexión. Una preparación procedente de la extensión mínima de
la sección~\ref{vii:vii:sec:corriente-extension-minima} puede determinar
el dato inicial del campo o del estado. Cada acción conserva su
variación propia: si se sustituye en una acción un campo minimizante
dependiente de la conexión, se aplica además su diferencial compuesto,
o el teorema de la envolvente bajo sus hipótesis de estacionariedad.
La acción~\eqref{vii:vii:eq:accion-dirac-material} realiza aquí un
acoplamiento afín en la contorsión.

\begin{proposition}[Corriente espinorial con normalización variacional]
\label{vii:vii:prop:corriente-dirac-explicita}
Sea
\(B^{IJK}=\bar\psi g^{[I}g^Jg^{K]}\psi\), donde los corchetes
designan la antisimetrización normalizada. La corriente definida en
\eqref{vii:vii:eq:corriente-variacional} por la acción
\eqref{vii:vii:eq:accion-dirac-material} es
\begin{equation}
 \sigma_{JK}
 =-\frac{\hbar}{2}\,
   \bar\psi\{g^I,\mathbb B_{JK}\}\psi\,\eta_I
 =-\frac{\hbar}{2}\,B^I{}_{JK}\eta_I.
 \label{vii:vii:eq:corriente-dirac-trivector}
\end{equation}
Es independiente de la contorsión cuando \(e\) y \(\psi\) están
fijos. La inversión de la sección~\ref{vii:vii:subsec:reconstruccion-torsion}
se aplica por tanto a esta corriente.
\end{proposition}
\begin{proof}
La variación de \(\Omega\) es
\(\delta\Omega=\frac12\delta\omega^{JK}\mathbb B_{JK}\).
Los dos sumandos cinéticos dan, respectivamente, el producto con
\(g^I\) a la izquierda y a la derecha. En consecuencia,
\[
 \delta_\omega S_D
 =\frac{\hbar}{4}\int_M\delta\omega^{JK}\wedge\eta_I\,
                 \bar\psi\{g^I,\mathbb B_{JK}\}\psi.
\]
La comparación con \eqref{vii:vii:eq:corriente-variacional} fija el
factor y el signo de \eqref{vii:vii:eq:corriente-dirac-trivector}.
Al expandir el anticonmutador y usar la relación de Clifford se
cancelan las contracciones de grado uno, y queda
\(\{g^I,\mathbb B_{JK}\}=g^{[I}g_Jg_{K]}\).
La acción es lineal en \(\Omega\), lo que prueba la última
afirmación.
\end{proof}

\subsection{Contracción efectiva y coeficiente del segundo momento}
\label{vii:vii:subsec:contraccion-espinorial}

La forma cuadrática lorentziana del trivector se escribe
\begin{equation}
 q(B)=\frac1{3!}B^{IJK}B_{IJK}
 =-(B^{012})^2-(B^{013})^2-(B^{023})^2+(B^{123})^2.
 \label{vii:vii:eq:cuadratica-trivector}
\end{equation}
Esta contracción conserva la firma del módulo espinorial. Su segundo
momento depende del estado y de la realización de los productos de
campos.

\begin{theorem}[Interacción efectiva de la corriente espinorial]
\label{vii:vii:thm:coeficiente-torsional-espinorial}
Con la acción material~\eqref{vii:vii:eq:accion-dirac-material}, la
eliminación algebraica de la contorsión da
\begin{equation}
 \mathcal L_{\rm spin}^{\rm eff}
 =C_\gamma q(B)\operatorname{vol}_e,\qquad
 C_\gamma=\frac{3\kappa c\hbar^2}{16}
                      \frac{\gamma^2}{1+\gamma^2}.
 \label{vii:vii:eq:coeficiente-cuadratico-espinorial}
\end{equation}
El coeficiente se obtiene de la corriente y de la inversa de Holst
con las normalizaciones de la acción precedente.
\end{theorem}
\begin{proof}
Por \eqref{vii:vii:eq:accion-efectiva-spin}, basta evaluar
\(-\frac14K_*^{IJ}\wedge\sigma_{IJ}\), con
\[
 Y^{IJ}=\kappa c\frac{\gamma^2}{1+\gamma^2}
           (-\star+\gamma^{-1}I)\sigma^{IJ},\qquad
 K_*=\mathsf C_e(Y).
\]
Aquí \(\mathsf C_e\) es la composición explícita de
\eqref{vii:vii:eq:torsion-explicita} con
\eqref{vii:vii:eq:contorsion-componentes}. La linealidad de esas dos
inversas reduce el cálculo a las partes \(-\star\) e \(I\).
Para dar la contracción completa, ordenamos las componentes del
trivector como \((012,013,023,123)\), extraemos los factores
\(\kappa c\hbar^2\gamma^2/(1+\gamma^2)\), y sustituimos
\(\sigma_{JK}=-\frac12B^I{}_{JK}\eta_I\). Las dos matrices
de la forma cuadrática restante son
\begin{equation}
 \mathcal B_{-\star}
 =\frac3{16}\operatorname{diag}(-1,-1,-1,1),\qquad
 \mathcal B_I=0_{4\times4}.
 \label{vii:vii:eq:matrices-contraccion-espinorial}
\end{equation}
Estas matrices se obtienen usando únicamente
\(e^I\wedge\eta_J=\delta_J^I\operatorname{vol}_e\): para
cada componente \(B^{abc}\), se forma \(\sigma\), se calcula
\(Y\), después
\(T^I=\sum_J\iota_{E_J}Y^{IJ}
 -\frac14e^I\wedge\sum_{L,J}\iota_{E_L}\iota_{E_J}Y^{LJ}\),
y finalmente
\(K_{IJK}=(T_{KIJ}-T_{IJK}-T_{JKI})/2\).
La contracción da las cuatro entradas diagonales mostradas; las
seis polarizaciones entre componentes distintas son cero para ambos
operadores. Es el cálculo de todos los coeficientes de las dos
formas cuadráticas. La parte \(\gamma^{-1}I\) se anula y la parte
\(-\star\) produce exactamente \eqref{vii:vii:eq:cuadratica-trivector},
con el coeficiente de \eqref{vii:vii:eq:coeficiente-cuadratico-espinorial}.
\end{proof}

\subsection{Realización fermiónica finita y segundo momento del estado}
\label{vii:vii:subsec:segundo-momento-car}

Fijamos una celda comóvil periódica, un marco ortonormal adaptado a su
observador y el modo espacial homogéneo de las cuatro componentes
espinoriales. Esta realización finita tiene espacio de estados
\(\mathcal F=\bigoplus_{n=0}^4\Lambda^n\mathbb C^4\).
Sean \(a_r,a_r^\dagger\), \(r=1,\ldots,4\), sus operadores
canónicos, con
\(\{a_r,a_s^\dagger\}=\delta_{rs}\) y
\(\{a_r,a_s\}=0\). El vacío de ocupación fija el orden
normal. Para una matriz \(M\), escribimos
\[
 \mathrm d\Gamma(M)=\sum_{r,s}a_r^\dagger M_{rs}a_s,
 \qquad \widehat N=\mathrm d\Gamma(I_4).
\]
La letra \(\Gamma\) de esta notación designa la segunda
cuantización; la holonomía nonádica conserva su notación propia.

Las cuatro matrices de los bilineales se calculan directamente:
\begin{equation}
 \begin{array}{c|cccc}
 abc&012&013&023&123\\ \hline
 M_{abc}=A_Dg^ag^bg^c
       &iJ\otimes R&iJ\otimes Z&iI_2\otimes J&iZ\otimes J
 \end{array}
 \label{vii:vii:eq:matrices-bilineales-espin}
\end{equation}
Todas son hermíticas, tienen traza cero y satisfacen
\(M_{abc}^2=I_4\). El observable cuártico correspondiente a
\eqref{vii:vii:eq:cuadratica-trivector}, en este orden normal declarado,
es
\begin{equation}
 \begin{split}
 \widehat{\mathscr Q}_{\rm sp}
 &=\sum_{a<b<c}\eta_{aa}\eta_{bb}\eta_{cc}
       :\!\mathrm d\Gamma(M_{abc})^2\!:\\
 &=\sum_{a<b<c}\eta_{aa}\eta_{bb}\eta_{cc}
       \mathrm d\Gamma(M_{abc})^2+2\widehat N.
 \end{split}
 \label{vii:vii:eq:observable-cuadratico-espin}
\end{equation}
El orden normal es parte de la realización finita especificada;
su definición hace explícita la sustracción de las contracciones de
una sola ocupación.

\begin{proposition}[Evaluación por las relaciones canónicas]
\label{vii:vii:prop:operador-espin-car}
El operador \(\widehat{\mathscr Q}_{\rm sp}\) es autoadjunto,
conserva cada sector de ocupación y tiene restricciones
\begin{equation}
 \widehat{\mathscr Q}_{\rm sp}|_{\Lambda^0\oplus\Lambda^1}=0,
 \qquad
 \widehat{\mathscr Q}_{\rm sp}|_{\Lambda^3}=4I_4,
 \qquad
 \widehat{\mathscr Q}_{\rm sp}|_{\Lambda^4}=8I_1.
 \label{vii:vii:eq:sectores-espin-car}
\end{equation}
En la base de \(\Lambda^2\mathbb C^4\) cuyos subconjuntos de
ocupación son \((12,13,23,14,24,34)\), su matriz es
\begin{equation}
 \begin{pmatrix}
 0&0&0&0&0&-4\\
 0&2&0&0&6&0\\
 0&0&2&-6&0&0\\
 0&0&-6&2&0&0\\
 0&6&0&0&2&0\\
 -4&0&0&0&0&0
 \end{pmatrix}.
 \label{vii:vii:eq:sector-dos-espin-car}
\end{equation}
\end{proposition}
\begin{proof}
La relación \(a_s a_t^\dagger=\delta_{st}-a_t^\dagger a_s\)
aplicada a \(\mathrm d\Gamma(M)^2\) separa su contracción de
una ocupación y da
\[
 :\!\mathrm d\Gamma(M)^2\!:
 =\mathrm d\Gamma(M)^2-\mathrm d\Gamma(M^2).
\]
Como los tres primeros signos de
\eqref{vii:vii:eq:matrices-bilineales-espin} son negativos y el último
positivo, la suma de las contracciones es \(-2\widehat N\).
Esto prueba \eqref{vii:vii:eq:observable-cuadratico-espin}.
La autoadjunción y la conservación de la ocupación se siguen de
esas mismas expresiones.

En \(\Lambda^0\) todos los sumandos se anulan. En \(\Lambda^1\),
\(\mathrm d\Gamma(M)=M\), y los cuadrados suman \(-2I_4\),
que se cancela con \(2\widehat N\). En \(\Lambda^4\),
\(\mathrm d\Gamma(M)=\operatorname{tr}M=0\), quedando ocho.
La identificación por complemento entre \(\Lambda^3\mathbb C^4\)
y el dual de \(\mathbb C^4\) lleva
\(\mathrm d\Gamma(M)\) a
\((\operatorname{tr}M)I-M^{\mathsf T}=-M^{\mathsf T}\).
Sus cuadrados vuelven a sumar \(-2I_4\), mientras
\(2\widehat N=6I_4\); por tanto la restricción es \(4I_4\).
Finalmente, la aplicación de las cuatro matrices explícitas a los
seis productos exteriores de la base indicada produce
\eqref{vii:vii:eq:sector-dos-espin-car}. Así queda evaluado el operador
en las dieciséis dimensiones de \(\mathcal F\).
\end{proof}

Para un operador de estado \(\varrho\geq0\),
\(\operatorname{Tr}\varrho=1\), definimos
\(q_\varrho=\operatorname{Tr}(\varrho\widehat{\mathscr Q}_{\rm sp})\).
Su dependencia de los segundos momentos es explícita:
\begin{equation}
 q_\varrho=\sum_{a<b<c}\eta_{aa}\eta_{bb}\eta_{cc}
 \left\{
  \langle\widehat B_{abc}\rangle_\varrho^2
  +\operatorname{Var}_\varrho(\widehat B_{abc})
  -\langle\widehat N\rangle_\varrho
 \right\},\qquad
 \widehat B_{abc}=\mathrm d\Gamma(M_{abc}).
 \label{vii:vii:eq:covarianza-espin-material}
\end{equation}
La covarianza y la sustracción normal se conservan incluso cuando
las corrientes medias se anulan. En particular, si \(P_3\) es la
proyección sobre \(\Lambda^3\mathbb C^4\),
\begin{equation}
 \varrho_3=\tfrac14P_3,\qquad
 \langle\widehat B_{abc}\rangle_{\varrho_3}=0,\qquad
 q_{\varrho_3}=4>0.
 \label{vii:vii:eq:testigo-positivo-espin}
\end{equation}
Es un estado positivo de traza uno y constituye un testigo explícito
del dominio positivo del funcional. El sector de dos ocupaciones
contiene también valores negativos; por ejemplo, el vector
\((|12\rangle+|34\rangle)/\sqrt2\) tiene valor \(-4\).
El signo de la intensidad procede así de la contracción y del
estado, y queda separado de la selección de condiciones cosmológicas.

\subsection{Densidad, presión y conservación de la amplitud comóvil}
\label{vii:vii:subsec:amplitud-espinorial-comovil}

En la realización espacialmente homogénea, fijamos
\[
 e^0=N(t)c\,\mathrm dt,\qquad e^j=a(t)\,\mathrm dx^j,
 \qquad \mathcal V(t)=a(t)^3V_c,
\]
con lapse \(N(t)>0\), factor de escala adimensional \(a(t)>0\)
y volumen comóvil \(V_c>0\). El marco y el modo periódico de la
subsección~\ref{vii:vii:subsec:segundo-momento-car} permanecen declarados.
La normalización del campo homogéneo es
\(\widehat\psi=\mathcal V^{-1/2}(a_1,a_2,a_3,a_4)^{\mathsf T}\).
Cada bilineal lleva un factor \(\mathcal V^{-1}\), y su producto
normal ordenado lleva \(\mathcal V^{-2}\). La forma temporal
cinética simétrica conserva esta normalización canónica: los términos
procedentes de derivar \(\mathcal V^{-1/2}\) se cancelan entre
sus dos sumandos.

\begin{theorem}[Amplitud torsional como funcional del estado]
\label{vii:vii:thm:amplitud-material-explicita}
Manteniendo fijos el estado canónico y la realización del producto
durante las variaciones métricas, el término efectivo de espín tiene
acción reducida
\begin{equation}
 S_{\rm spin}^{\rm hom}
 =\int\!\mathrm dt\,N(t)
       \frac{cC_\gamma}{a(t)^3V_c}\,q_\varrho.
 \label{vii:vii:eq:accion-homogenea-espin}
\end{equation}
Su densidad de energía y su presión isotrópica son
\begin{equation}
 \rho_{\rm spin}=p_{\rm spin}
 =-\frac{cC_\gamma}{V_c^2}\,q_\varrho a^{-6}.
 \label{vii:vii:eq:densidad-presion-espin}
\end{equation}
En cualquier evolución que conserve \(q_\varrho\), la amplitud
constante correspondiente es
\begin{equation}
 \boxed{\displaystyle
 s_0[\varrho,V_c]
 =\frac{3\kappa c^2\hbar^2}{16}
    \frac{\gamma^2}{1+\gamma^2}
    \frac{\operatorname{Tr}
       (\varrho\widehat{\mathscr Q}_{\rm sp})}{V_c^2}.}
 \label{vii:vii:eq:amplitud-s0-espinorial}
\end{equation}
\end{theorem}
\begin{proof}
La integración espacial de
\eqref{vii:vii:eq:coeficiente-cuadratico-espinorial} multiplica el
factor \(\mathcal V^{-2}\) por
\(Nc\mathcal V\,\mathrm dt\), lo que da
\eqref{vii:vii:eq:accion-homogenea-espin}. En variables canónicas,
la variación del lapse define la energía mediante
\(\delta_N S_{\rm m}=-\int E\,\delta N\,\mathrm dt\).
Por tanto,
\[
 E_{\rm spin}=-\frac{cC_\gamma q_\varrho}{a^3V_c},
 \qquad \rho_{\rm spin}=\frac{E_{\rm spin}}{a^3V_c}.
\]
La variación espacial isotrópica satisface
\(\delta_a S_{\rm m}=\int3Na^2V_c p\,\delta a\,\mathrm dt\).
Aplicada a \eqref{vii:vii:eq:accion-homogenea-espin}, da
\[
 3Na^2V_c p_{\rm spin}
 =-\frac{3NcC_\gamma q_\varrho}{a^4V_c}.
\]
Esto prueba simultáneamente el signo, la presión y la densidad de
\eqref{vii:vii:eq:densidad-presion-espin}. Al conservarse
\(q_\varrho\), el coeficiente de \(-a^{-6}\) es constante y
se obtiene \eqref{vii:vii:eq:amplitud-s0-espinorial}. La dimensionalidad
es consistente: \([cC_\gamma]=\text{energía}\,L^3\),
de modo que \(s_0\) tiene dimensión de densidad de energía.
\end{proof}

\begin{corollary}[Dominio conservado de amplitud positiva]
\label{vii:vii:cor:dominio-positivo-conservado}
Sea una evolución unitaria de \(\mathcal F\) que preserve
\(\Lambda^3\mathbb C^4\). Para todo estado inicial de traza uno
soportado en ese sector, se cumple \(q_{\varrho(t)}=4\), y
\begin{equation}
 s_0=\frac{3\kappa c^2\hbar^2}{4V_c^2}
                   \frac{\gamma^2}{1+\gamma^2}>0.
 \label{vii:vii:eq:amplitud-positiva-sector-tres}
\end{equation}
\end{corollary}
\begin{proof}
La invariancia del sector conserva el soporte del estado; en él,
\eqref{vii:vii:eq:sectores-espin-car} da
\(\widehat{\mathscr Q}_{\rm sp}=4I\). La conservación de la
traza fija entonces su esperanza en cuatro, independientemente de
la evolución interna de sus coherencias. La positividad de los
acoplamientos, \(V_c>0\) y \(\gamma\ne0\) prueba la fórmula
y su signo.
\end{proof}

Para estados generales, la conservación de
\(\langle\widehat N\rangle\) determina la dilución de la densidad
de ocupación, \(n(t)=\langle\widehat N\rangle/(a^3V_c)\).
La conservación del momento cuártico se comprueba por separado.
Si \(i\hbar\dot\varrho=[H(t),\varrho]\), entonces
\begin{equation}
 \frac{\mathrm d q_\varrho}{\mathrm dt}
 =\frac{i}{\hbar}\operatorname{Tr}
       \bigl(\varrho[H(t),\widehat{\mathscr Q}_{\rm sp}]\bigr).
 \label{vii:vii:eq:conservacion-momento-cuartico}
\end{equation}
La conmutación, la invariancia de un sector donde el observable es
escalar, o una conservación demostrada para el estado concreto
proporcionan las condiciones pertinentes. Si el momento evoluciona,
\eqref{vii:vii:eq:densidad-presion-espin} conserva su evaluación
instantánea; su balance incluye
\[
 \dot\rho_{\rm spin}+6H\rho_{\rm spin}
 =-\frac{cC_\gamma}{V_c^2}\,a^{-6}\dot q_\varrho,
 \qquad H=\dot a/a
\]
en tiempo propio. Este término registra el intercambio con los
restantes grados de libertad materiales.

La composición queda determinada: representación espinorial del
transporte, acción material, corriente variacional, inversión de
Cartan--Holst, segundo momento y densidad normalizada. El dominio
\(s_0>0\) de la sección~\ref{vii:vii:sec:dilucion-rebote} dispone así
de un funcional evaluable y de estados testigo explícitos. Su uso en
la dinámica cosmológica conserva además las condiciones materiales,
geométricas y de conservación enunciadas en aquella sección.
